% lint-skip: 未完成ステータス（figure/table caption 未対応，人手チェック時に解消予定）
\documentclass[../../main.tex]{subfiles}
\begin{document}

\begin{flushright}
  \footnotesize\textit{【自動文字起こし・要確認】}
\end{flushright}

{\bf [解答]} 原点Oを支点とし，点Xの位置ベクトル$\vec{x}$とする．

\[
  \begin{cases}
    \vec{a}_1 = \dfrac{1}{3}(\vec{a}+2\vec{b}) \\
    \vec{b}_1 = \dfrac{1}{3}(\vec{b}+2\vec{c}) \\
    \vec{c}_1 = \dfrac{1}{3}(\vec{c}+2\vec{a})
  \end{cases}
\]
\[
  \begin{cases}
    \vec{a}_2 = \dfrac{1}{3}(\vec{a}_1+2\vec{b}_1) \\
    \vec{b}_2 = \dfrac{1}{3}(\vec{b}_1+2\vec{c}_1) \\
    \vec{c}_2 = \dfrac{1}{3}(\vec{c}_1+2\vec{a}_1)
  \end{cases}
\]
から
\[
  \begin{cases}
    \vec{a}_2 = \dfrac{1}{9}(\vec{a}+4\vec{b}+4\vec{c}) \\
    \vec{b}_2 = \dfrac{1}{9}(4\vec{a}+\vec{b}+4\vec{c}) \\
    \vec{c}_2 = \dfrac{1}{9}(4\vec{a}+4\vec{b}+\vec{c})
  \end{cases}
\]
より，
\begin{align}
  |A_2B_2| &= |\vec{b}_2-\vec{a}_2| = \dfrac{1}{3}|\vec{a}-\vec{b}| = \dfrac{1}{3}|AB| \\
  |B_2C_2| &= |\vec{c}_2-\vec{b}_2| = \dfrac{1}{3}|\vec{b}-\vec{c}| = \dfrac{1}{3}|BC| \\
  |C_2A_2| &= |\vec{a}_2-\vec{c}_2| = \dfrac{1}{3}|\vec{c}-\vec{a}| = \dfrac{1}{3}|CA|
\end{align}
だから三辺相等より
\[ \triangle ABC \sim \triangle A_2B_2C_2 \quad // \]

\begin{figure}[htbp]
  \centering
  \begin{tikzpicture}[scale=1.65,line join=round]
    \coordinate (A) at (1,2.25);
    \coordinate (B) at (0,0);
    \coordinate (C) at (3.2,0);
    \coordinate (A1) at ($(A)!2/3!(B)$);
    \coordinate (B1) at ($(B)!2/3!(C)$);
    \coordinate (C1) at ($(C)!2/3!(A)$);
    \coordinate (A2) at ($(A1)!2/3!(B1)$);
    \coordinate (B2) at ($(B1)!2/3!(C1)$);
    \coordinate (C2) at ($(C1)!2/3!(A1)$);
    \fill[orange!10] (A2)--(B2)--(C2)--cycle;
    \draw[blue!75!black,very thick] (A)--(B)--(C)--cycle;
    \draw[teal!70!black,thick] (A1)--(B1)--(C1)--cycle;
    \draw[orange!85!black,very thick] (A2)--(B2)--(C2)--cycle;
    \foreach \name/\pos in {A/above,B/below left,C/below right}
    \node[\pos] at (\name) {$\name$};
    \node[left,teal!70!black] at (A1) {$A_1$};
    \node[below,teal!70!black] at (B1) {$B_1$};
    \node[right,teal!70!black] at (C1) {$C_1$};
    \node[below,orange!85!black] at (A2) {$A_2$};
    \node[right,orange!85!black] at (B2) {$B_2$};
    \node[above,orange!85!black] at (C2) {$C_2$};
    \node[left,font=\scriptsize] at ($(A)!0.5!(A1)$) {$2$};
    \node[left,font=\scriptsize] at ($(A1)!0.5!(B)$) {$1$};
    \node[below,font=\scriptsize] at ($(B)!0.5!(B1)$) {$2$};
    \node[below,font=\scriptsize] at ($(B1)!0.5!(C)$) {$1$};
    \node[right,font=\scriptsize] at ($(C)!0.5!(C1)$) {$2$};
    \node[right,font=\scriptsize] at ($(C1)!0.5!(A)$) {$1$};
  \end{tikzpicture}
  \caption{各辺を $2:1$ に内分してできる $A_1B_1C_1$ と $A_2B_2C_2$}
\end{figure}

\end{document}
